\chapter{拟射影簇}\label{chap:quasiprojective-varieties}

\noindent\textbf{第 \thechapter{} 章学习目标}
\begin{itemize}
  \item 定义并运用开集, 闭集, 闭包, 稠密性和不可约性等拓扑概念.
  \item 定义拟射影簇并描述其例子, 包括仿射的, 射影的以及既非仿射也非射影的例子.
  \item 描述拟射影簇情形下的正则映射和正则函数, 理解它们的连续性和局部性.
  \item 理解拟射影簇的局部仿射性质.
  \item 定义拟射影簇的积, 并判断积的子集是否闭.
  \item 理解为什么射影簇在正则映射下的像是闭的.
\end{itemize}

我们已经研究了两类簇: 仿射簇和射影簇. 它们在某些方面表现相似, 却似乎属于两个不同的世界. 本章的目标是引入一种更一般的\textquotedbl{}簇\textquotedbl{}的概念, 把仿射簇和射影簇都作为特例, 从而统一这两个世界. 关键的认识是, 通过包含映射 $j_{0}:\mathbb{A}^{n}\to\mathbb{P}^{n}$, 任意仿射簇 $X$ 都可视为其射影闭包 $\overline{X}$ 的子集. 由此产生的射影簇的子集在拓扑上是特殊的: 它们在 $\overline{X}$ 上的 Zariski 拓扑中开, 我们将在本章第一节精确地定义这一拓扑. 另一方面, 整个射影簇 $\overline{X}$ 在其自身的 Zariski 拓扑中也是开集, 所以 $\mathbb{P}^{n}$ 中自然有一类子集, 同时包含仿射簇和射影簇, 即射影簇的开子集. 我们把这些称为\emph{拟射影簇}.

研究拟射影簇时, 有几个方面需要格外留意. 我们需要引入比此前更多的拓扑背景, 并且需要确保为拟射影簇建立的每个概念 (例如\textquotedbl{}正则映射\textquotedbl{}的定义) 在仿射情形和射影情形中都专化为相应的概念, 只要该概念已经在相应情形下有定义. 不过, 这项工作的一项回报是, 我们将能够为射影簇定义许多迄今尚未建立的概念. 例如, 射影簇的维数是什么? 如何判断射影簇是否光滑? 一旦不仅能研究簇本身, 还能够研究其任一开子集, 而其中有许多是本书已经充分研究过的仿射簇, 我们就能够把这些概念引入射影理论.

\par\smallskip
{\footnotesize
\noindent\textcopyright{} The Author(s) 2025\par
\noindent E. Clader, D. Ross, \emph{Beginning in Algebraic Geometry}, Undergraduate Texts in Mathematics,
\href{https://doi.org/10.1007/978-3-031-88819-9_12}{\nolinkurl{https://doi.org/10.1007/978-3-031-88819-9_12}}.\par}
